\documentclass[11pt, a4paper]{article}

% Compatibility for pdfLaTeX (Overleaf default)
\usepackage[utf8]{inputenc}
\usepackage[T1]{fontenc}
\usepackage{textcomp}

% Support for the Rupee symbol (₹) in pdfLaTeX
\usepackage{tfrupee}
\DeclareUnicodeCharacter{20B9}{\rupee}

\usepackage{geometry}
\geometry{margin=1in}
\usepackage{tabularx}
\usepackage{booktabs}
\usepackage{hyperref}
\usepackage{enumitem}
\usepackage{xcolor}
\usepackage{titlesec}

% Link styling
\hypersetup{
    colorlinks=true,
    linkcolor=blue,
    filecolor=magenta,      
    urlcolor=cyan,
}

\title{\textbf{Exhaustive Codebase Analysis \& Engineering Breakdown}\\UPI Offline Mesh}
\author{}
\date{}

\begin{document}

\maketitle

\hrule
\vspace{1em}

\section*{A. Executive Summary}
The \textbf{UPI Offline Mesh} project is a Java/Spring Boot backend demonstration of a \textbf{Mesh-Routed Deferred Settlement System}. It solves the problem of peer-to-peer financial transactions in zero-connectivity environments (like basements or rural areas). 

It operates by having the sender's device create a cryptographically secure, time-bound payload that hops across untrusted intermediate devices via a simulated Bluetooth mesh (Gossip Protocol). When any device in the mesh gains internet access (acting as a ``Bridge Node''), it uploads the encrypted payload to this Spring Boot backend. The backend strictly enforces exactly-once processing (Idempotency) via atomic hashing, decrypts the payload using Hybrid Encryption (RSA + AES-GCM), verifies freshness to prevent replay attacks, and safely settles the transaction in an ACID-compliant database. 

It proves that decentralized, offline payments can be securely routed through untrusted networks and safely reconciled centrally without double-spending.

\section*{B. Folder Structure Analysis}

\begin{verbatim}
UPI_Without_Internet-main/
├── .mvn/wrapper/                  # Maven Wrapper files.
├── src/
│   ├── main/
│   │   ├── java/com/demo/upimesh/ # Root Java package.
│   │   │   ├── config/            # Configuration classes (e.g., background scheduling).
│   │   │   ├── controller/        # REST APIs and MVC controllers.
│   │   │   ├── crypto/            # Cryptography logic (Key generation, Hybrid Crypto).
│   │   │   ├── model/             # JPA Entities (DB tables) and DTOs.
│   │   │   ├── service/           # Core business logic (Idempotency, Settlement, Mesh).
│   │   │   └── UpiMeshApplication.java # Spring Boot main entry point.
│   │   └── resources/
│   │       ├── application.properties # App configuration (Port, DB dialect, TTLs).
│   │       └── templates/             # HTML templates (Thymeleaf, Dashboard UI).
│   └── test/
│       └── java/com/demo/upimesh/ # JUnit/Spring Boot test files.
├── mvnw & mvnw.cmd                # Unix/Windows execution scripts for Maven Wrapper.
├── pom.xml                        # Project Object Model (dependencies, plugins).
└── README.md                      # Extensive project documentation.
\end{verbatim}

\section*{C. Technology Stack Breakdown}

\begin{enumerate}
    \item \textbf{Spring Boot (v3.3.5)}: The core application framework.
    \begin{itemize}
        \item \textit{Why}: Provides auto-configuration, embedded Tomcat, and rapid REST API development.
    \end{itemize}
    \item \textbf{Spring Data JPA \& Hibernate}: The Object-Relational Mapping (ORM) layer.
    \begin{itemize}
        \item \textit{Why}: Abstracts raw SQL into Java objects (\texttt{@Entity}). Handles \texttt{@Transactional} states and \texttt{@Version} for optimistic locking.
    \end{itemize}
    \item \textbf{H2 Database}: An in-memory SQL database.
    \begin{itemize}
        \item \textit{Why}: Perfect for demos. It spins up and tears down with the app, requiring zero local setup from the user.
    \end{itemize}
    \item \textbf{Thymeleaf}: A server-side Java template engine.
    \begin{itemize}
        \item \textit{Why}: Used to serve \texttt{dashboard.html}. Allows the backend to easily render the UI on \texttt{localhost:8080} without needing a separate React/Node frontend.
    \end{itemize}
    \item \textbf{Java Cryptography Architecture (JCA)}: Core Java libraries (\texttt{javax.crypto}).
    \begin{itemize}
        \item \textit{Why}: Used for RSA and AES algorithms, secure random number generation, and SHA-256 hashing without external dependencies like BouncyCastle.
    \end{itemize}
    \item \textbf{Jackson (\texttt{com.fasterxml.jackson})}: JSON serialization/deserialization library.
    \begin{itemize}
        \item \textit{Why}: Converts Java objects (\texttt{PaymentInstruction}) into byte arrays for encryption, and parses incoming HTTP JSON requests.
    \end{itemize}
    \item \textbf{Tailwind CSS \& Lucide Icons} (Frontend via CDN):
    \begin{itemize}
        \item \textit{Why}: Used in the \texttt{dashboard.html} for rapid, modern UI styling without a build step.
    \end{itemize}
\end{enumerate}

\section*{D. Architecture Overview}

\subsection*{End-to-End Workflow}
\begin{enumerate}
    \item \textbf{Creation (Simulated Offline)}: The \texttt{DemoService} acts as the sender's phone. It creates a \texttt{PaymentInstruction} containing VPAs, amount, PIN, a UUID nonce, and a timestamp.
    \item \textbf{Encryption (Offline)}: The sender encrypts the payload using the Server's Public RSA Key. Specifically, it generates a throwaway AES-256 key, encrypts the JSON with AES-GCM, encrypts the AES key with RSA-OAEP, and packs them into a base64 \texttt{MeshPacket}.
    \item \textbf{Gossip (Offline)}: The sender hands the packet to their \texttt{VirtualDevice}. The \texttt{MeshSimulatorService} periodically copies packets between devices, decrementing the TTL.
    \item \textbf{Uplink (Online)}: When a device with \texttt{hasInternet=true} holds the packet, the \texttt{/api/mesh/flush} endpoint triggers it to POST the packet to the server's \texttt{/api/bridge/ingest} endpoint.
    \item \textbf{Ingestion \& Deduplication (Server)}:
    \begin{itemize}
        \item \texttt{BridgeIngestionService} receives the POST.
        \item It hashes the ciphertext (\texttt{SHA-256}).
        \item It attempts to claim the hash in \texttt{IdempotencyService} (using \texttt{putIfAbsent}).
        \item If claimed, it proceeds. If not (another thread/bridge delivered it), it drops it as a duplicate.
    \end{itemize}
    \item \textbf{Decryption \& Validation (Server)}: The server uses its Private RSA Key to unwrap the AES key, then decrypts the payload. It checks the GCM auth tag for tampering and the \texttt{signedAt} timestamp against the freshness window (24h).
    \item \textbf{Settlement (Server)}: \texttt{SettlementService} opens a DB \texttt{@Transactional}. It checks balances, verifies the PIN hash, debits the sender, credits the receiver, and writes an immutable \texttt{Transaction} record.
\end{enumerate}

\section*{E. Feature-by-Feature Analysis}

\subsection*{1. Hybrid Cryptography Engine}
\begin{itemize}
    \item \textbf{Purpose}: Secure data over untrusted intermediaries.
    \item \textbf{Implementation}: \texttt{HybridCryptoService.java}. RSA cannot encrypt large payloads. It generates a symmetric AES-256 key, uses it to encrypt the JSON payload in GCM mode (which provides an Authentication Tag to detect tampering), and then uses RSA-2048 to encrypt only the AES key.
\end{itemize}

\subsection*{2. Atomic Idempotency Cache}
\begin{itemize}
    \item \textbf{Purpose}: Prevent double-spending when multiple bridge nodes upload the same transaction.
    \item \textbf{Implementation}: \texttt{IdempotencyService.java}. Wraps a \texttt{ConcurrentHashMap}. The \texttt{claim()} method uses \texttt{putIfAbsent()}. This is a JVM-level atomic operation. Even if three threads hit it at the exact same nanosecond, only one succeeds.
\end{itemize}

\subsection*{3. Bluetooth Mesh Simulator}
\begin{itemize}
    \item \textbf{Purpose}: Replicate offline peer-to-peer data transfer without physical hardware.
    \item \textbf{Implementation}: \texttt{MeshSimulatorService.java}. Holds an in-memory map of \texttt{VirtualDevice} objects. The \texttt{gossipOnce()} method iterates through all devices and copies packets to other devices if the packet's TTL is $>$ 0. It runs on a Spring \texttt{@Scheduled} loop every 5 seconds.
\end{itemize}

\subsection*{4. ACID Ledger Settlement}
\begin{itemize}
    \item \textbf{Purpose}: Safely move money and record the event.
    \item \textbf{Implementation}: \texttt{SettlementService.java}. Uses \texttt{@Transactional}. If \texttt{accounts.save()} fails, the whole method rolls back. Uses \texttt{@Version} on the \texttt{Account} entity so Hibernate enforces Optimistic Locking at the SQL level (e.g., \texttt{UPDATE accounts SET balance = X, version = 2 WHERE vpa = Y AND version = 1}).
\end{itemize}

\section*{F. File-by-File Analysis}

\subsection*{\texttt{/crypto}}
\begin{itemize}
    \item \textbf{\texttt{ServerKeyHolder.java}}: \texttt{@PostConstruct} creates an RSA-2048 KeyPair on boot. Provides getters for the private/public keys. \textit{Role}: Simulates an HSM (Hardware Security Module).
    \item \textbf{\texttt{HybridCryptoService.java}}: Implements \texttt{encrypt()}, \texttt{decrypt()}, and \texttt{hashCiphertext()}. \textit{Role}: The cryptographic core.
\end{itemize}

\subsection*{\texttt{/model} (Domain Layer)}
\begin{itemize}
    \item \textbf{\texttt{Account.java}}: JPA \texttt{@Entity} for user balances. Has \texttt{@Version} for concurrency control.
    \item \textbf{\texttt{Transaction.java}}: JPA \texttt{@Entity} for the immutable ledger. Note the \texttt{@Index(unique = true)} on \texttt{packetHash} as a database-level defense against duplicates.
    \item \textbf{\texttt{MeshPacket.java}}: DTO representing the wire format. Contains public routing info (\texttt{ttl}, \texttt{packetId}) and the opaque \texttt{ciphertext}.
    \item \textbf{\texttt{PaymentInstruction.java}}: The actual cleartext business logic (VPAs, amount, nonce, time). Only exists before encryption and after decryption.
    \item \textbf{\texttt{*Repository.java}}: Spring Data interfaces generating SQL queries dynamically.
\end{itemize}

\subsection*{\texttt{/service} (Business Logic)}
\begin{itemize}
    \item \textbf{\texttt{IdempotencyService.java}}: The deduplication gatekeeper. Evicts old hashes via \texttt{@Scheduled}.
    \item \textbf{\texttt{BridgeIngestionService.java}}: The pipeline orchestrator. Hashes $\rightarrow$ Claims $\rightarrow$ Decrypts $\rightarrow$ Checks Time $\rightarrow$ Calls Settlement. Returns early if any step fails.
    \item \textbf{\texttt{SettlementService.java}}: The financial engine. Executes DB logic.
    \item \textbf{\texttt{DemoService.java}}: Seeds DB with accounts on boot. Contains \texttt{createPacket()} which simulates the Android client building and encrypting a payload.
    \item \textbf{\texttt{MeshSimulatorService.java} / \texttt{VirtualDevice.java}}: The network simulation layer handling the gossip protocol.
\end{itemize}

\section*{G. Concept and Topic Breakdown}

\begin{table}[h]
\centering
\begin{tabularx}{\textwidth}{|>{\bfseries}l|X|X|}
\hline
Concept & What it is & Why it's used here \\ \hline
Gossip Protocol & Epidemic network routing where nodes share data with neighbors. & To pass packets offline until they reach the internet. \\ \hline
Hybrid Encryption & Combining Asymmetric (RSA) and Symmetric (AES) encryption. & RSA is too slow/limited for large JSON; AES is fast but needs a secure key exchange. \\ \hline
AEAD (AES-GCM) & Authenticated Encryption. Encrypts data and provides a cryptographic checksum (Tag). & To ensure untrusted phones cannot alter the payment amount or VPAs. \\ \hline
Idempotency & A property where an operation applied multiple times yields the same result as applied once. & To prevent the ``Duplicate Storm'' of multiple bridges delivering the same payment. \\ \hline
Compare-and-Set & An atomic processor instruction (or data structure operation). & \texttt{putIfAbsent} guarantees thread-safe deduplication without blocking locks. \\ \hline
Optimistic Locking & DB locking technique using version numbers instead of row locks. & Prevents ``Lost Updates'' during concurrent DB writes without performance penalties. \\ \hline
Nonce & ``Number used once''. A UUID inside the payload. & Ensures two identical ₹100 payments create completely different ciphertexts. \\ \hline
\end{tabularx}
\end{table}

\section*{H. Security and Scalability Analysis}

\subsection*{Security Considerations}
\begin{itemize}
    \item \textbf{Strengths}:
    \begin{itemize}
        \item \textbf{Tamper-Proof}: AES-GCM ensures any altered bit throws an exception.
        \item \textbf{E2EE}: Intermediaries only see base64 gibberish.
        \item \textbf{Replay Protection}: The \texttt{signedAt} field and 24h freshness window prevent attackers from saving a ciphertext and submitting it a year later.
    \end{itemize}
    \item \textbf{Risks \& Technical Debt (Acknowledged in README)}:
    \begin{itemize}
        \item \textbf{Offline Double Spend}: Because validation is deferred, a user with ₹500 can send ₹500 to Alice offline, then walk away and send ₹500 to Bob offline. Whichever packet reaches the server first settles; the second gets \texttt{REJECTED} (NSF). The receiver assumes the risk.
        \item \textbf{Server Key Storage}: Currently generated in memory on boot. In prod, must be in an HSM/AWS KMS.
        \item \textbf{No Bridge Auth}: \texttt{/api/bridge/ingest} is open. Needs mTLS or JWTs.
    \end{itemize}
\end{itemize}

\subsection*{Scalability Considerations}
\begin{itemize}
    \item \textbf{Current Bottleneck}: \texttt{ConcurrentHashMap} for idempotency is strictly tied to a single JVM memory space.
    \item \textbf{Production Path}: To scale to multiple backend instances, the \texttt{ConcurrentHashMap} MUST be replaced with \textbf{Redis} using the \texttt{SET key value NX EX 86400} command. This ensures idempotency is distributed across the cluster.
    \item \textbf{Database}: H2 must be replaced with PostgreSQL.
\end{itemize}

\section*{I. Rebuild-from-Scratch Roadmap}

If you want to build this exact architecture from scratch, follow this learning order:

\begin{enumerate}
    \item \textbf{Java \& Spring Boot Basics}:
    \begin{itemize}
        \item Learn to set up a Spring Boot project, define \texttt{@RestController}, and handle \texttt{@RequestBody}.
    \end{itemize}
    \item \textbf{Concurrency in Java (Crucial)}:
    \begin{itemize}
        \item Understand Race Conditions.
        \item Learn \texttt{ConcurrentHashMap}, \texttt{AtomicInteger}, \texttt{CountDownLatch}, and \texttt{ExecutorService}. Build a simple dedup script.
    \end{itemize}
    \item \textbf{Relational DBs \& Hibernate (JPA)}:
    \begin{itemize}
        \item Learn \texttt{@Entity}, \texttt{@Repository}, and how to use the H2 in-memory DB.
        \item \textit{Advanced}: Study Database Transactions (\texttt{@Transactional}) and Optimistic Locking (\texttt{@Version}).
    \end{itemize}
    \item \textbf{Applied Cryptography in Java}:
    \begin{itemize}
        \item Learn how to use \texttt{java.security.MessageDigest} (SHA-256).
        \item Learn \texttt{javax.crypto.Cipher}. Implement basic AES encryption.
        \item Learn RSA KeyPair generation.
        \item \textit{Mastery}: Combine them into a Hybrid Encryption wrapper using \texttt{AES/GCM/NoPadding} and \texttt{RSA/ECB/OAEPWithSHA-256AndMGF1Padding}.
    \end{itemize}
    \item \textbf{Distributed Systems Concepts}:
    \begin{itemize}
        \item Study Idempotency keys.
        \item Study the Gossip Protocol and Time-To-Live (TTL) networks.
    \end{itemize}
    \item \textbf{Testing}:
    \begin{itemize}
        \item Learn JUnit 5. Write tests that intentionally spawn multiple threads to break your code, then fix it.
    \end{itemize}
    \item \textbf{Frontend Integration}:
    \begin{itemize}
        \item Learn basic Thymeleaf to serve HTML from Spring, and vanilla JS \texttt{fetch()} to interact with your APIs.
    \end{itemize}
\end{enumerate}

\end{document}
